\section{\textsc{ProprioVLA}}

\textsc{ProprioVLA} adds a proprioceptive grounding path to a VLA policy while
leaving the standard action-supervised objective intact. During training, the
method mines self-grounding targets from state-aligned residual motion and uses
them to supervise learned grounding queries over visual features. At test time,
the mined heatmaps are not required.

\begin{figure}[t]
    \centering
    \includegraphics[width=\linewidth]{propriovla_main_figure_preview.png}
    \caption{
    Overview of \textsc{ProprioVLA}. Given an observation, instruction, and
    proprioceptive state, the base VLA produces a standard multimodal
    representation. \textsc{ProprioVLA} adds a lightweight learned-query
    grounding adapter over visual tokens and supervises its attention with
    mined self-grounding heatmaps. The heatmap supervision is used only during
    training.
    }
    \label{fig:propriovla_method}
\end{figure}

\subsection{Self-Grounding from State-Aligned Motion}

The robot reveals itself through motion that is synchronized with its state
change. Given consecutive observations \((o_t,o_{t+1})\) and proprioceptive
states \((s_t,s_{t+1})\), we define
\(\Delta s_t=s_{t+1}-s_t\). We estimate dense visual motion \(F_t(p)\) between
frames and subtract global or background motion to obtain residual motion
\(R_t(p)\) at image location \(p\). This residual suppresses camera drift and
static-scene variation, emphasizing motion caused by the robot or by objects
directly coupled to the robot.

We convert residual motion into a confidence-weighted self-grounding target.
Let \(\mathrm{sal}(R_t(p))\) denote local residual-motion saliency and let
\[
    \alpha_t(p)
    =
    \left[
    \cos\bigl(\psi_R(R_t(p)),\psi_s(\Delta s_t)\bigr)
    \right]_+
\]
measure state-motion alignment after lightweight projections \(\psi_R\) and
\(\psi_s\). The unnormalized heatmap score is
\[
    e_t(p)=\mathrm{sal}(R_t(p))\,\alpha_t(p).
\]
We normalize \(e_t\) over the image or patch grid to obtain \(H_t\), and set
a confidence \(c_t\in[0,1]\) from the total aligned-motion mass. Low confidence
downweights frames where the self-motion cue is weak or ambiguous. Thus the
training target is mined from interaction alone, without manual masks, external
grounding datasets, or stronger foundation-model supervision.

\subsection{Learned Queries for Self-Grounding}

In \textsc{ProprioVLA}, a small set of learned queries forms the grounding path
that localizes the robot in the image. Let
\(Q=\{q_m\}_{m=1}^{M}\) be \(M\) learned query latents, and let
\(K_t,V_t=f_v(o_t)\) be visual keys and values from the VLA vision encoder.
The queries attend over visual patches according to
\[
    A_t(m,p)
    =
    \mathrm{softmax}_{p}
    \left(
    \frac{\bar{q}_m \bar{K}_t(p)^\top}{\tau}
    \right),
    \qquad
    A_t(p)=\frac{1}{M}\sum_{m=1}^{M}A_t(m,p),
\]
where \(\bar{q}_m\) and \(\bar{K}_t(p)\) are \(\ell_2\)-normalized query and
key projections, and \(\tau\) is a temperature. The averaged attention map
\(A_t\) is the model's predicted self-grounding distribution over visual
patches.

\subsection{Normalized KL Grounding Objective}

The mined heatmap supervises the learned-query attention through a normalized
KL objective. Let \(U_t\) denote the uniform distribution over valid visual
tokens. The grounding loss is
\[
    \mathcal{L}_{\mathrm{ground}}
    =
    c_t
    \frac{
    \mathrm{KL}(H_t\,\|\,A_t)
    }{
    \mathrm{KL}(H_t\,\|\,U_t)+\epsilon
    } ,
\]
after resizing \(H_t\) to the attention grid and normalizing it over valid
tokens. The denominator makes the auxiliary loss comparable across target
sharpness and token counts: matching the target approaches zero, while uniform
attention is normalized to approximately one. The complete training loss is
\[
    \mathcal{L}
    =
    \mathcal{L}_{\mathrm{act}}
    +
    \lambda_g \mathcal{L}_{\mathrm{ground}} .
\]
Here \(\mathcal{L}_{\mathrm{act}}\) is the standard imitation loss used to
fine-tune the VLA.
